\chapter{Using GitHub for Your Coursework}
\label{app:github}

The \prg platform lives on GitHub, and so should your work \citep{prg32repo}. This
appendix is a practical introduction to the version-control workflow you will use
to obtain the platform, track your own changes, recover from mistakes, and submit
assignments. It assumes no prior experience with Git. The aim is not Git mastery
but the handful of operations that make coursework safer and submission painless.

\section{Why version control, concretely}

Two everyday disasters motivate everything in this appendix. The first is the
embedding lab of Chapter~\ref{ch:ccap}, where you temporarily modify the firmware's
build list and main loop, then must restore them; without a record of the original,
``restore'' becomes guesswork. The second is the moment, familiar to every
programmer, when a working game stops working after an edit and you cannot remember
what you changed. Version control solves both: it remembers every committed state of
your files and lets you compare or return to any of them.

\begin{prgconcept}
\textbf{Git} records snapshots, called \emph{commits}, of your project over time.
Each commit is a labelled, recoverable state. \textbf{GitHub} hosts your Git
repository online, so your work is backed up, shareable, and submittable. Together
they turn ``I think I broke something'' into ``let me see exactly what I changed.''
\end{prgconcept}

\section{Getting the platform}

You obtain \prg by \emph{cloning} its repository, which downloads the full project
and its history \citep{prg32repo}.

\lstinputlisting[
    style=bash,
    caption={%\fname{clone\_repo.sh} \protect\
       Cloning the platform repository.}
]{scripts/app_d_github/clone_repo.sh}

To incorporate later updates from the upstream project without losing your own
work, you fetch and merge changes; for coursework, however, the cleaner path is to
keep your own changes on a branch, described below.

\section{The everyday cycle}

Three commands cover the great majority of daily use: \code{status} to see what has
changed, \code{add} to stage the changes you want to keep, and \code{commit} to
record them with a message.

\lstinputlisting[
    style=bash,
    caption={%\fname{git\_cycle.sh} \protect\
       The everyday Git cycle: inspect, stage, commit.}
]{scripts/app_d_github/git_cycle.sh}

\begin{prgwarn}
Write commit messages that your future self will understand: ``Add boundary clamp
to mygame paddle,'' not ``stuff'' or ``fix.'' A week later, when you are hunting for
the commit where a bug appeared, descriptive messages are the difference between a
minute and an hour. The platform's own history follows this practice.
\end{prgwarn}

\section{Branches: experiment freely}

A \emph{branch} is an independent line of work. Before a risky change, such as the
firmware embedding of Chapter~\ref{ch:ccap}, or a large refactor of your game, make
a branch. If the experiment fails, you discard the branch and your main line is
untouched; if it succeeds, you merge it back.

\lstinputlisting[
    style=bash,
    caption={%\fname{branch\_workflow.sh} \protect\
       Working on a branch, then merging or discarding.}
]{scripts/app_d_github/branch_workflow.sh}

\begin{prgconcept}
A branch lets you try something that might break your project, in isolation. The
embedding lab is the textbook case: branch first, embed and test, then either merge
the working result or discard the branch to restore the firmware to its original
state in one command: exactly the ``restore afterwards'' the tutorial requires
\citep{prg32tutorialc}.
\end{prgconcept}

\section{Recovering from mistakes}

When an edit breaks a working game, Git shows you precisely what changed since your
last commit, and can undo it.

\lstinputlisting[
    style=bash,
    caption={%\fname{undo\_changes.sh} \protect\
       Seeing and undoing uncommitted changes.}
]{scripts/app_d_github/undo_changes.sh}

The discipline that makes this powerful is committing often: after each working
step of the incremental method this book teaches. If you commit after every
verified change (``draw the paddle,'' ``add input,'' ``add the clamp''), then any
break is, at worst, one step's worth of work to inspect or undo.

\section{Submitting your work}

For submission you push your commits to your own GitHub repository, then share its
link or open a pull request as your course directs. You create an empty repository
on GitHub, connect it as a remote, and push.

\lstinputlisting[
    style=bash,
    caption={%\fname{publish\_repo.sh} \protect\
       Publishing your work to your own GitHub repository.}
]{scripts/app_d_github/publish_repo.sh}

\begin{prgaside}
The \prg repository models the conventions you are learning: a \fname{README} that
orients a newcomer, a \fname{CONTRIBUTING} guide, a \fname{LICENSE}, a
\fname{CITATION.cff} so the work can be cited, and continuous-integration workflows
that run automated checks on every change \citep{prg32repo}. When you structure your
own coursework repository, these are the files that signal care and make your work
easy for an instructor to evaluate.
\end{prgaside}

\section{Citing the platform}

If you use \prg in a report, a thesis, or any submitted work, cite it. The
repository ships citation metadata in \fname{CITATION.cff} for exactly this purpose
\citep{prg32repo}. This book cites the platform throughout, and the bibliography
entry for the repository is the model to follow in your own writing.

\begin{prgcheck}
You should now be able to clone the platform, see what you have changed, stage and
commit changes with a meaningful message, work safely on a branch, undo an
unwanted change, and push your work to GitHub for submission. These few operations
are enough for an entire first year of coursework.
\end{prgcheck}
